\section{Introduction}
\label{sec:Introduction}

The Epsom Oddballs Monopoly Run is an example of the Orienteering Problem. In this essay we discuss how optimal routes can be found, use this analysis to make observations on strategies for planning a route as a human, and briefly comment on how such a game should be designed. The main focus will be on the 2026 board, but data from 2025 is sometimes used as a comparison. Places of the board\footnote{Where a square can be achieved by visiting multiple places, a variety of places are given that attempt to capture where a team may go. For squares that could be achieved in lots of different ways this is sometimes implemented as a single place, such as ``Any coffee shop", ``Any Chinese restaurant", and ``Any kebab shop" where there are instances that are extremely near to Metropolis.} are given in table~\ref{tab:Places} and all code and data is available at~\cite{GithubRepository}. The rules are as follows.

\begin{enumerate}
	\item Stay within the 70 minute time limit.
	\item Get points for each square you visit, with more points typically allocated to further away places.
	\item A square is worth double points if you have visited all squares within a group.
\end{enumerate}

There are bonus goals and points that can be obtained from other teams visiting squares after other teams, but since these rely on what other teams do, these aspects are considered out of scope. The goal instead will be to find routes that achieve the maxium number of points only from visiting places.

Section~\ref{sec:FormulatingTheProblem} details how the problem is solved, with further details about solving problems on a graph with the chosen approach given in section~\ref{sec:NetworkConstraints}. The results are given in section~\ref{sec:Results}.

% Do not edit - this file was automatically generated by output_places

\begin{longtable}{ccccc}
\caption{The 2025 Monopoly Run places with their associated squares, groups, and values.} \label{tab:Places} \\
\toprule
Place ID & Place & Square & Value & Group Name \\
\midrule
\endfirsthead
\caption[]{The 2025 Monopoly Run places with their associated squares, groups, and values.} \\
\toprule
Place ID & Place & Square & Value & Group Name \\
\midrule
\endhead
\midrule
\multicolumn{5}{r}{Continued on next page} \\
\midrule
\endfoot
\bottomrule
\endlastfoot
M & Metropolis & Metropolis & 0 & Terminals \\
1 & Glyn School & Glyn School & 5 & School \\
2 & Nescott College & Nescott College & 6 & School \\
3 & Danetree Primary School & Danetree Primary School & 6 & School \\
4 & Ewell East Train Station & Ewell East Train Station & 10 & Station \\
5 & Nova Forca & Nova Forca & 5 & Run \\
6 & Nonsuch & Parkrun Start Nonsuch & 6 & Run \\
7 & Harrier Centre & Running Track & 6 & Run \\
8 & Cheam Recreation Ground & Running Track & 6 & Run \\
9 & Hobbledown Farm & Hobbledown Farm & 7 & Horton \\
10 & Adventure Golf & Adventure Golf & 6 & Horton \\
11 & Long Grove Skate Park & Long Grove Skate Park & 4 & Horton \\
12 & Stoneleigh Train Station & Stoneleigh Train Station & 10 & Station \\
13 & The Watch House & The Watch House & 7 & Cheam \\
14 & Nonsuch Mansion House & Nonsuch Mansion House & 10 & Cheam \\
15 & Ewell Court House & Ewell Court House & 10 & Cheam \\
16 & Epsom Playhouse & Epsom Playhouse & 3 & Halls \\
17 & Epsom Odeon & Epsom Odeon & 2 & Halls \\
18 & Opposite Roseberry Park & UCA & 3 & Halls \\
19 & Opposite Pets at Home & UCA & 3 & Halls \\
20 & Near Back of Odeon & UCA & 3 & Halls \\
21 & Epsom Train Station & Epsom Train Station & 10 & Station \\
22 & Epsom Town Hall & Epsom Town Hall & 2 & Landmarks \\
23 & The Clock Tower & The Clock Tower & 2 & Landmarks \\
24 & Epsom Racecourse & Epsom Racecourse & 10 & Landmarks \\
25 & Pizzerium / Pizza Express & Pizza Restuarant & 2 & Other \\
26 & Papa John / Ask Italian & Pizza Restuarant & 2 & Other \\
27 & Dominos (Epsom) & Pizza Restuarant & 2 & Other \\
28 & Fireaway / American Village & Pizza Restuarant & 2 & Other \\
29 & Village Pizza (Ewell Bypass) & Pizza Restuarant & 2 & Other \\
30 & Tops Pizza (Chessington Road) & Pizza Restuarant & 2 & Other \\
31 & Pizza Express (Cheam) & Pizza Restuarant & 2 & Other \\
32 & Ntunetta (Ewell) & Pizza Restuarant & 2 & Other \\
33 & Epsom Air Scouts & Scout Hut & 4 & Other \\
34 & Epsom Methodist Scout Group & Scout Hut & 4 & Other \\
35 & Ewell Rainsters Scout Group & Scout Hut & 4 & Other \\
36 & St Martins Scout Group & Scout Hut & 4 & Other \\
37 & Little Acres Lodge & Little Acres Lodge & 8 & Other \\
38 & Ewell West Train Station & Ewell West Train Station & 10 & Station \\
39 & Epsom Hospital & Epsom Hospital & 4 & Fluids \\
40 & Jolly Cooper Pub & Jolly Cooper Pub & 4 & Fluids \\
41 & Epsom Well & Epsom Well & 7 & Fluids \\
\end{longtable}
